\setcounter{section}{4}


  \zwischenueberschrift{Irreduzible affin-algebraische Mengen}


\inputdefinition


{}


{


Eine
\definitionsverweis {affin-algebraische Menge}{}{}


\mavergleichskette


{\vergleichskette


{V
}


{ \subseteq }{ { {\mathbb A}_{  K }^{  n   } }
}


{  }{
}


{    }{
}


{   }{
}


}
{}{}{}
heißt
\definitionswort {irreduzibel}{,}
wenn


\mavergleichskette


{\vergleichskette


{V
}


{ \neq }{ \emptyset
}


{  }{
}


{    }{
}


{   }{
}


}
{}{}{}
ist und es keine Zerlegung


\mavergleichskette


{\vergleichskette


{V
}


{ = }{ Y \cup Z
}


{  }{
}


{    }{
}


{   }{
}


}
{}{}{}
mit
\definitionsverweis {affin-algebraischen Mengen}{}{}


\mavergleichskette


{\vergleichskette


{Y,Z
}


{ \subset }{ V
}


{  }{
}


{    }{
}


{   }{
}


}
{}{}{}
gibt.


}


Die Zariski-abgeschlossene Menge $V$ ist also irreduzibel genau dann, wenn


\mavergleichskettedisp


{\vergleichskette


{ V
}


{ \neq} { \emptyset
}


{ } {
}


{ } {
}


{ } {
}


}
{}{}{}
ist und eine Zerlegung


\mavergleichskette


{\vergleichskette


{ V
}


{ = }{ Y \cup Z
}


{  }{
}


{    }{
}


{   }{
}


}
{}{}{}
nur mit


\mavergleichskette


{\vergleichskette


{ V
}


{ = }{ Y
}


{  }{
}


{    }{
}


{   }{
}


}
{}{}{}
oder mit


\mavergleichskette


{\vergleichskette


{ V
}


{ = }{ Z
}


{  }{
}


{    }{
}


{   }{
}


}
{}{}{}
möglich ist. Dasselbe folgt dann sofort für endliche Darstellungen.


Die Irreduzibilität ist eine rein topologische Eigenschaft, wobei man die obige Definition mit abgeschlossenen Mengen formulieren muss anstatt mit affin-algebraischen Mengen
\zusatzklammer {den abgeschlossenen Mengen in der Zariski-Topologie} {} {.}


Die folgenden Bilder zeigen einige (nicht) irreduzible affin-algebraische Teilmengen. Was sind dabei die irreduziblen Komponenten
\zusatzklammer {siehe unten} {} {?}


\bild{ \begin{center}


\includegraphics[width=5.5cm]{ \bildeinlesungpng {Delaunay_points} {png} }


\end{center}


\bildtext {} }


\bildlizenz { Delaunay points.png } {} {Nü Es} {Commons} {CC-BY-SA-3.0} {}


\bild{ \begin{center}


\includegraphics[width=5.5cm]{ \bildeinlesungsvg {Gerade} {svg} }


\end{center}


\bildtext {} }


\bildlizenz { Gerade.svg } {} {Lokilech} {Commons} {CC-BY-SA-3.0} {}


\bild{ \begin{center}


\includegraphics[width=5.5cm]{ \bildeinlesungsvg {Straight_lines} {svg} }


\end{center}


\bildtext {} }


\bildlizenz { Straight lines.svg } {} {Lokilech} {Commons} {CC-BY-SA-3.0} {}


\bild{ \begin{center}


\includegraphics[width=5.5cm]{ \bildeinlesungpng {Linear_space2} {png} }


\end{center}


\bildtext {} }


\bildlizenz { Linear space2.png } {} {Kmhkmh} {Commons} {CC-BY-SA-3.0} {}


\inputbeispiel{}


{


Wir betrachten den affinen Raum


\mathl{{ {\mathbb A}_{  K }^{  n   } }}{.} Wenn $K$ endlich ist, so besteht der Raum nur aus endlich vielen Punkten und nur die einpunktigen Teilmengen sind irreduzibel. Insbesondere ist der affine Raum außer bei


\mavergleichskette


{\vergleichskette


{ n
}


{ = }{ 0
}


{  }{
}


{    }{
}


{   }{
}


}
{}{}{}
nicht irreduzibel.


Bei unendlichem $K$ ist der affine Raum


\mathl{{ {\mathbb A}_{  K }^{  n   } }}{} hingegen irreduzibel. Es sei nämlich


\mavergleichskette


{\vergleichskette


{ { {\mathbb A}_{  K }^{  n   } }
}


{ = }{ Y \cup Z
}


{  }{
}


{    }{
}


{   }{
}


}
{}{}{}
mit echt kleineren affin-algebraischen Teilmengen. D.h. für die offenen Komplemente


\mavergleichskette


{\vergleichskette


{ U
}


{ = }{ { {\mathbb A}_{  K }^{  n   } }  \setminus Y
}


{  }{
}


{    }{
}


{   }{
}


}
{}{}{}
und


\mavergleichskette


{\vergleichskette


{ W
}


{ = }{ { {\mathbb A}_{  K }^{  n   } } \setminus Z
}


{  }{
}


{    }{
}


{   }{
}


}
{}{}{}
ist einerseits


\mavergleichskette


{\vergleichskette


{ U,W
}


{ \neq }{ \emptyset
}


{  }{
}


{    }{
}


{   }{
}


}
{}{}{,}
aber


\mavergleichskette


{\vergleichskette


{ U \cap W
}


{ = }{ \emptyset
}


{  }{
}


{    }{
}


{   }{
}


}
{}{}{.}
Das widerspricht aber

Aufgabe 3.20
.


}


\inputfaktbeweis


{Affine Varietäten/Irreduzible Teilmengen/Primideale/Fakt}


{Lemma}


{}


{


\faktsituation {}


\faktvoraussetzung {Es sei


\mavergleichskette


{\vergleichskette


{V
}


{ \subseteq }{ { {\mathbb A}_{  K }^{  n   } }
}


{  }{
}


{    }{
}


{   }{
}


}
{}{}{}
eine
\definitionsverweis {affin-algebraische Menge}{}{}
mit
\definitionsverweis {Verschwindungsideal}{}{}


\mathl{\operatorname{Id} \,(V)}{.}}


\faktfolgerung {Dann ist $V$ genau dann
\definitionsverweis {irreduzibel}{}{,}
wenn


\mathl{\operatorname{Id} \, (V)}{} ein
\definitionsverweis {Primideal}{}{}
ist.}


\faktzusatz {}


\faktzusatz {}


}


{


\teilbeweis {}{}{}


{Es sei


\mathl{\operatorname{Id} \, (V)}{} kein Primideal. Bei


\mavergleichskette


{\vergleichskette


{ \operatorname{Id} \, (V)
}


{ = }{ K[X_1 , \ldots , X_n]
}


{  }{
}


{    }{
}


{   }{
}


}
{}{}{}
ist


\mavergleichskette


{\vergleichskette


{ V
}


{ = }{ \emptyset
}


{  }{
}


{    }{
}


{   }{
}


}
{}{}{,}
also ist $V$ nicht irreduzibel nach Definition. Andernfalls gibt es Polynome


\mavergleichskette


{\vergleichskette


{ F,G
}


{ \in }{ K[X_1 , \ldots ,  X_n]
}


{  }{
}


{    }{
}


{   }{
}


}
{}{}{}
mit


\mavergleichskette


{\vergleichskette


{ F G
}


{ \in }{ \operatorname{Id} \, (V)
}


{  }{
}


{    }{
}


{   }{
}


}
{}{}{,}
aber


\mavergleichskette


{\vergleichskette


{ F,G
}


{ \notin }{ \operatorname{Id}\,(V)
}


{  }{
}


{    }{
}


{   }{
}


}
{}{}{.}
Dies bedeutet, dass es Punkte


\mavergleichskette


{\vergleichskette


{ P,Q
}


{ \in }{ V
}


{  }{
}


{    }{
}


{   }{
}


}
{}{}{}
mit


\mavergleichskette


{\vergleichskette


{ F(P)
}


{ \neq }{ 0
}


{  }{
}


{    }{
}


{   }{
}


}
{}{}{}
und


\mavergleichskette


{\vergleichskette


{ G(Q)
}


{ \neq }{ 0
}


{  }{
}


{    }{
}


{   }{
}


}
{}{}{}
gibt. Wir betrachten die beiden Ideale


\mavergleichskette


{\vergleichskette


{ {\mathfrak a}_1
}


{ = }{ \operatorname{Id}\, (V)+(F)
}


{  }{
}


{    }{
}


{   }{
}


}
{}{}{}
und


\mavergleichskette


{\vergleichskette


{ {\mathfrak a}_2
}


{ = }{ \operatorname{Id} \, (V)+(G)
}


{  }{
}


{    }{
}


{   }{
}


}
{}{}{.}
Daher ist


\mavergleichskettedisp


{\vergleichskette


{ V( {\mathfrak a}_1 ), V( {\mathfrak a}_2 )
}


{ \subseteq} { V(\operatorname{Id}\, (V)  )
}


{ =} { V
}


{ } {
}


{ } {
}


}
{}{}{}
nach

Lemma 3.8

 (3)
.
Wegen
\mathkor {} {P \not\in V({\mathfrak a}_1)} {und} {Q \not\in V({\mathfrak a}_2)} {}
sind diese Inklusionen echt. Andererseits ist


\mavergleichskettedisp


{\vergleichskette


{ V({\mathfrak a}_1) \cup V({\mathfrak a}_2)
}


{ =} { V({\mathfrak a}_1 \cdot {\mathfrak a}_2)
}


{ =} { V( \operatorname{Id} \,(V)  )
}


{ =} { V
}


{ } {
}


}
{}{}{,}
sodass eine nicht-triviale Zerlegung von $V$ vorliegt und somit $V$ nicht irreduzibel ist.}


{}
\teilbeweis {}{}{}


{Es sei nun $V$ nicht irreduzibel. Bei


\mavergleichskette


{\vergleichskette


{ V
}


{ = }{ \emptyset
}


{  }{
}


{    }{
}


{   }{
}


}
{}{}{}
ist


\mavergleichskette


{\vergleichskette


{ \operatorname{Id}\, (V)
}


{ = }{ K[X_1 , \ldots , X_n]
}


{  }{
}


{    }{
}


{   }{
}


}
{}{}{}
kein Primideal. Es sei also


\mavergleichskette


{\vergleichskette


{ V
}


{ \neq }{ \emptyset
}


{  }{
}


{    }{
}


{   }{
}


}
{}{}{}
mit der nicht-trivialen Zerlegung


\mavergleichskette


{\vergleichskette


{ V
}


{ = }{ Y \cup Z
}


{  }{
}


{    }{
}


{   }{
}


}
{}{}{.}
Es sei


\mavergleichskette


{\vergleichskette


{ Y
}


{ = }{ V \left(  {\mathfrak a}_1   \right)
}


{  }{
}


{    }{
}


{   }{
}


}
{}{}{}
und


\mavergleichskette


{\vergleichskette


{ Z
}


{ = }{ V \left(  {\mathfrak a}_2  \right)
}


{  }{
}


{    }{
}


{   }{
}


}
{}{}{.}
Wegen


\mavergleichskette


{\vergleichskette


{ Y
}


{ \subset }{ V
}


{  }{
}


{    }{
}


{   }{
}


}
{}{}{}
gibt es einen Punkt


\mathbed {P \in V = V(\operatorname{Id}\, (V))} {}


 {P \not\in V( {\mathfrak a}_1 )} {}


 {} {} {} {.}
Also gibt es auch ein


\mavergleichskette


{\vergleichskette


{ F
}


{ \in }{ {\mathfrak a}_1
}


{  }{
}


{    }{
}


{   }{
}


}
{}{}{,}


\mavergleichskette


{\vergleichskette


{ F(P)
}


{ \neq }{ 0
}


{  }{
}


{    }{
}


{   }{}


}
{}{}{,}
und somit


\mavergleichskette


{\vergleichskette


{ F
}


{ \notin }{ \operatorname{Id} \, (V)
}


{  }{
}


{    }{
}


{   }{
}


}
{}{}{.}
Ebenso gibt es


\mavergleichskette


{\vergleichskette


{ G
}


{ \in }{ {\mathfrak a}_2
}


{  }{
}


{    }{
}


{   }{
}


}
{}{}{,}


\mavergleichskette


{\vergleichskette


{ G
}


{ \notin }{ \operatorname{Id} \, (V)
}


{  }{
}


{    }{
}


{   }{}


}
{}{}{.}
Für einen beliebigen Punkt


\mavergleichskette


{\vergleichskette


{ Q
}


{ \in }{ V
}


{ = }{ Y \cup Z
}


{    }{
}


{   }{
}


}
{}{}{}
ist


\mavergleichskette


{\vergleichskette


{ (FG)(Q)
}


{ = }{ 0
}


{  }{
}


{    }{
}


{   }{
}


}
{}{}{,}
da $F$ auf $Y$ und $G$ auf $Z$ verschwindet. Also ist


\mavergleichskette


{\vergleichskette


{ FG
}


{ \in }{ \operatorname{Id} \, (V)
}


{  }{
}


{    }{
}


{   }{
}


}
{}{}{}
und daher ist


\mathl{\operatorname{Id} \, (V)}{} kein Primideal.}


{}


}


\inputdefinition


{}


{


Es sei $V$ eine
\definitionsverweis {affin-algebraische Menge}{}{.}
Eine affin-algebra\-ische Teilmenge


\mavergleichskette


{\vergleichskette


{W
}


{ \subseteq }{ V
}


{  }{
}


{    }{
}


{   }{
}


}
{}{}{}
heißt eine \definitionswort {irreduzible Komponente}{} von $V$, wenn sie
\definitionsverweis {irreduzibel}{}{}
ist und wenn es keine irreduzible Teilmenge


\mavergleichskette


{\vergleichskette


{W
}


{ \subset }{ W'
}


{ \subseteq }{ V
}


{    }{
}


{   }{
}


}
{}{}{}
gibt.


}


Ist $V$ irreduzibel, so ist $V$ selbst die einzige irreduzible Komponente von $V$. Wir werden in

Satz 9.11

sehen, dass jede affin-algebraische Menge sich als eine endliche Vereinigung von irreduziblen Komponenten schreiben lässt.


\inputbeispiel{}


{


Wir betrachten die Gleichung


\mavergleichskettedisp


{\vergleichskette


{ F
}


{ =} { Y^2+X^2(X+1)^2
}


{ =} { 0
}


{ } {
}


{ } {
}


}
{}{}{.}
In den reellen Zahlen hat diese Gleichung zwei Lösungen: da ein reelles Quadrat nie negativ ist, kann $F$ nur dann $0$ sein, wenn beide Summanden null sind, und das impliziert einerseits


\mavergleichskette


{\vergleichskette


{ Y
}


{ = }{ 0
}


{  }{
}


{    }{
}


{   }{
}


}
{}{}{}
und andererseits


\mavergleichskette


{\vergleichskette


{ X
}


{ = }{ 0
}


{  }{
}


{    }{
}


{   }{
}


}
{}{}{}
oder


\mavergleichskette


{\vergleichskette


{ X
}


{ = }{ -1
}


{  }{
}


{    }{
}


{   }{
}


}
{}{}{.}
Insbesondere ist die reelle Lösungsmenge nicht zusammenhängend und nicht irreduzibel
\zusatzklammer {und das Verschwindungsideal zur reellen Situation ist sehr groß} {} {.}


Betrachtet man $F$ dagegen über den komplexen Zahlen, so gibt es eine Faktorisierung


\mavergleichskettedisp


{\vergleichskette


{F
}


{ =} { (Y+ { \mathrm i} X(X+1))(Y- { \mathrm i} X(X+1))
}


{ } {
}


{ } {
}


{ } {
}


}
{}{}{}
in irreduzible Polynome. Dies zeigt zugleich, dass $F$ als Polynom in


\mathl{\R [X,Y]}{} irreduzibel ist
\zusatzklammer {obwohl das reelle Nullstellengebilde nicht irreduzibel ist} {} {.}
Die Nullstellenmenge über den komplexen Zahlen besteht aus den beiden Graphen


\mavergleichskette


{\vergleichskette


{ Y
}


{ = }{ \pm { \mathrm i} X(X+1)
}


{  }{
}


{    }{
}


{   }{
}


}
{}{}{,}
die sich in
\mathkor {} {(0,0)} {und} {(-1,0)} {}
schneiden.


Bei der Gleichung


\mathl{Y^2+Z^2+X^2(X+1)^2}{} gibt es wieder nur zwei reelle Lösungspunkte, das Polynom ist aber sowohl reell als auch komplex betrachtet irreduzibel.


}


\bild{ \begin{center}


\includegraphics[width=5.5cm]{ \bildeinlesungjpg {Hydrant_Insel_Krk_Kroatien} {jpg} }


\end{center}


\bildtext {} }


\bildlizenz { Hydrant Insel Krk Kroatien.jpg } {} {Usien} {Commons} {CC-BY-SA-3.0} {}


\inputbeispiel{}


{


Wir betrachten im affinen Raum


\mathl{{ {\mathbb A}_{  K  }^{  3   } }}{}
\zusatzklammer {


\mavergleichskettek


{\vergleichskettek


{ K
}


{ = }{ \R
}


{  }{
}


{    }{
}


{   }{
}


}
{}{}{}} {} {}
die beiden \stichwort {Zylinder} {}


\mathdisp {S_1= \left\{  (x,y,z)   \mid   x^2+y^2 = 1         \right\}  \text{ und }  S_2= \left\{  (x,y,z)   \mid   y^2+z^2 = 1        \right\}} { . }


Das sind beides irreduzible Mengen, wie wir später sehen werden
\zusatzklammer {für $K$ unendlich} {} {.}
Wie sieht ihr Durchschnitt aus? Der Durchschnitt wird durch das Ideal ${\mathfrak a}$ beschrieben, das durch


\mathl{X^2+Y^2-1}{} und


\mathl{Y^2+Z^2-1}{} erzeugt wird. Zieht man die eine Gleichung von der anderen ab, so erhält man


\mavergleichskettedisp


{\vergleichskette


{ X^2-Z^2
}


{ =} { (X-Z)(X+Z)
}


{ \in} { {\mathfrak a}
}


{ } {
}


{ } {
}


}
{}{}{.}
Die beiden einzelnen Faktoren gehören aber nicht zu ${\mathfrak a}$, da beispielsweise


\mathl{(1,0,-1)}{} ein Punkt des Schnittes ist, an dem


\mathl{X-Z}{} nicht verschwindet
\zusatzklammer {Charakteristik $\neq 2$} {} {,}
und


\mathl{(1,0,1)}{} ein Punkt des Schnittes ist, an dem


\mathl{X+Z}{} nicht verschwindet. Die Komponenten des Schnittes werden vielmehr durch


\mathdisp {{\mathfrak b}_1= {\mathfrak a} + (X-Z) \text{ und } {\mathfrak b}_2= {\mathfrak a} + (X+Z)} {  }


beschrieben. Das sind beides Primideale, der Restklassenring ist


\mavergleichskettedisp


{\vergleichskette


{ K[X,Y,Z]/({\mathfrak b}_1)
}


{ =} { K[X,Y,Z]/({\mathfrak a} + (X-Z))
}


{ =} { K[X,Y]/ \left(  X^2+Y^2-1  \right)
}


{ } {
}


{ } {
}


}
{}{}{.}
Um die zweite Gleichung einzusehen, eliminiert man $Z$ mit der hinteren Gleichung, und die beiden Zylindergleichungen werden dann identisch. Ebenso ist die Argumentation für das andere Ideal. Geometrisch gesprochen heißt dies, dass ein Punkt des Durchschnittes


\mathl{S_1 \cap S_2}{} in der Ebene


\mavergleichskette


{\vergleichskette


{ E_1
}


{ = }{ V(Z-X)
}


{  }{
}


{    }{
}


{   }{
}


}
{}{}{}
oder in der Ebene


\mavergleichskette


{\vergleichskette


{ E_2
}


{ = }{ V(Z+X)
}


{  }{
}


{    }{
}


{   }{
}


}
{}{}{}
liegt. Es ist


\mavergleichskettedisp


{\vergleichskette


{ E_1 \cap S_1
}


{ =} { E_1 \cap S_1 \cap S_2
}


{ =} { E_1 \cap S_2
}


{ } {
}


{ } {
}


}
{}{}{}
und ebenso für $E_1$, da auf diesen Ebenen die beiden Zylindergleichungen identisch werden.


Wie sehen die Durchschnitte in den Ebenen aus? Wir betrachten die Ebene $E_1$ mit den Koordinaten $Y$ und


\mavergleichskette


{\vergleichskette


{ U
}


{ = }{ Z+X
}


{  }{
}


{    }{
}


{   }{
}


}
{}{}{.}
Es ist dann


\mavergleichskette


{\vergleichskette


{ X
}


{ = }{ \frac{1}{2} ((Z+X) - (Z-X))
}


{  }{
}


{    }{
}


{   }{
}


}
{}{}{}
und damit kann man die erste Zylindergleichung als


\mavergleichskettedisp


{\vergleichskette


{ \left(  \frac{1}{2} ((Z+X) - (Z-X))  \right)^2+Y^2
}


{ =} { 1
}


{ } {
}


{ } {
}


{ } {
}


}
{}{}{}
schreiben. Auf der Ebene $E_1$, die ja durch


\mavergleichskette


{\vergleichskette


{ Z
}


{ = }{ X
}


{  }{
}


{    }{
}


{   }{
}


}
{}{}{}
festgelegt ist, wird aus dieser Gleichung


\mavergleichskettedisp


{\vergleichskette


{ \left(  \frac{1}{2} U  \right)^2+Y^2
}


{ =} { 1
}


{ } {
}


{ } {
}


{ } {
}


}
{}{}{}
also


\mavergleichskette


{\vergleichskette


{ \frac{1}{4}U^2+Y^2
}


{ = }{ 1
}


{  }{
}


{    }{
}


{   }{
}


}
{}{}{.}
Dies ist die Gleichung einer \stichwort {Ellipse} {,} was auch anschaulich klar ist. Man beachte, dass in der obigen Berechnung des Restklassenringes


\mathl{K[X,Y,Z]/({\mathfrak b}_1)}{} aber eine Kreisgleichung auftritt. Dies sollte deshalb nicht überraschen, da Kreis und Ellipse durch eine lineare Variablentransformation ineinander überführbar sind und dass daher insbesondere die Restklassenringe isomorph sind. Als \stichwort {metrisches Gebilde} {} sind Kreis und Ellipse verschieden, und der Durchschnitt der beiden Zylinder besteht aus zwei Ellipsen. Bei einer \stichwort {orthonormalen Variablentransformation} {} bleibt die metrische Struktur erhalten. Die Variablen


\mathl{Y, \left(  X+Z  \right) , \left(  X-Z  \right)}{} definieren aber keine orthonormale Transformation.


Halten wir also fest: Der Durchschnitt der beiden Zylinder ist


\mavergleichskettedisp


{\vergleichskette


{ S_1 \cap S_2
}


{ =} { V({\mathfrak b}_1) \cup V({\mathfrak b}_2)
}


{ } {
}


{ } {
}


{ } {
}


}
{}{}{,}
wobei


\mavergleichskette


{\vergleichskette


{ {\mathfrak b}_1
}


{ = }{ (X^2+Y^2-1, X-Z )
}


{  }{
}


{    }{
}


{   }{
}


}
{}{}{}
und


\mavergleichskette


{\vergleichskette


{ {\mathfrak b}_2
}


{ = }{ (X^2+Y^2-1, X+Z )
}


{  }{
}


{    }{
}


{   }{
}


}
{}{}{}
zwei Ellipsen beschreiben.


Wie liegen diese beiden Ellipsen zueinander? Dazu berechnen wir ihren Durchschnitt, der durch die Summe von ${\mathfrak b}_1$ und ${\mathfrak b}_2$ beschrieben wird. Es ist


\mavergleichskettedisp


{\vergleichskette


{ {\mathfrak b}_1 + {\mathfrak b}_2
}


{ =} { \left(  X^2+Y^2-1 , X-Z, X+Z  \right)
}


{ =} { \left(  Y^2-1,X,Z  \right)
}


{ } {
}


{ } {
}


}
{}{}{.}
Die Lösungsmenge davon besteht aus den beiden Punkten
\mathkor {} {(0,1,0)} {und} {(0,-1,0)} {.}


}


\bild{ \begin{center}


\includegraphics[width=5.5cm]{ \bildeinlesungsvg {Cylinder_principal_directions} {svg} }


\end{center}


\bildtext {} }


\bildlizenz { Cylinder principal directions.svg } {Luca Antonelli} {Luke Antony} {Commons} {CC-BY-SA-3.0} {}


  \zwischenueberschrift{Zur Anzahl der Punkte auf Kurven}


Wir haben bereits gesehen, dass der Schnitt einer Kurve mit einer Geraden nur aus endlich vielen Punkten besteht, es sei denn, die Gerade ist selbst eine Komponente der Kurve
\zusatzklammer {siehe

Lemma 1.3
} {} {.}
Dies wollen wir zunächst auf den Schnitt von zwei beliebigen ebenen Kurven verallgemeinern. Als Hilfsmittel benötigen wir die folgende Definition.


\inputdefinition


{}


{


Es sei $K$ ein
\definitionsverweis {Körper}{}{}
und


\mathl{K[X]}{} der
\definitionsverweis {Polynomring in einer Variablen}{}{}
über $K$. Dann nennt man den
\definitionsverweis {Quotientenkörper}{}{}


\mathl{Q(K[X])}{} den \definitionswort {rationalen Funktionenkörper}{} über $K$
\zusatzklammer {oder \definitionswort {Körper der rationalen Funktionen}{} über $K$} {} {.}
Er wird mit


\mathl{K(X)}{} bezeichnet.


}


\inputfaktbeweis


{Ebene Kurven/Schnitt ohne Komponenten/Endlich viele Punkte/Fakt}


{Satz}


{}


{


\faktsituation {}


\faktvoraussetzung {Es sei $K$ ein
\definitionsverweis {Körper}{}{}
und seien


\mavergleichskette


{\vergleichskette


{F,G
}


{ \in }{ K[X,Y]
}


{  }{
}


{    }{
}


{   }{
}


}
{}{}{}
zwei Polynome ohne gemeinsamen nichtkonstanten Faktor.}


\faktfolgerung {Dann gibt es nur endlich viele Punkte


\mathl{P_1 , \ldots , P_n}{} mit


\mavergleichskette


{\vergleichskette


{P_i
}


{ \in }{ V(F,G)
}


{  }{
}


{    }{
}


{   }{
}


}
{}{}{.}}


\faktzusatz {}


\faktzusatz {}


}


{


Wir betrachten


\mavergleichskette


{\vergleichskette


{ F,G
}


{ \in }{ K[X,Y]
}


{  }{
}


{    }{
}


{   }{
}


}
{}{}{}
als Elemente in


\mathl{K(X)[Y]}{,} wobei


\mathl{K(X)}{} den
\definitionsverweis {Körper der rationalen Funktionen}{}{}
in $X$ bezeichne. Es haben dann nach

Aufgabe 4.27

auch $F$ und $G$ keinen gemeinsamen Teiler in


\mathl{K(X)[Y]}{.} Da dieser Ring ein
\definitionsverweis {Hauptidealbereich}{}{}
ist, erzeugen sie zusammen das Einheitsideal, d.h. es gibt Polynome


\mavergleichskette


{\vergleichskette


{ A,B
}


{ \in }{ K(X)[Y]
}


{  }{
}


{    }{
}


{   }{
}


}
{}{}{}
mit


\mavergleichskette


{\vergleichskette


{ AF+BG
}


{ = }{ 1
}


{  }{
}


{    }{
}


{   }{
}


}
{}{}{.}
Multiplikation mit dem Hauptnenner von $A$ und $B$ ergibt in


\mathl{K[X,Y]}{} die Gleichung


\mavergleichskette


{\vergleichskette


{ \tilde{A} F + \tilde{B}G
}


{ = }{ H
}


{  }{
}


{    }{
}


{   }{
}


}
{}{}{}
mit


\mavergleichskette


{\vergleichskette


{ H
}


{ \in }{ K[X]
}


{  }{
}


{    }{
}


{   }{
}


}
{}{}{.}
Eine gemeinsame Nullstelle in


\mathl{{\mathbb A}^{2}_{ K }}{} von $F$ und von $G$ muss also eine Nullstelle von $H$ sein. Es gibt also nur endlich viele Werte für $X$, für die eine gemeinsame Nullstelle vorliegt. Wenn man die Rollen von $X$ und von $Y$ vertauscht, so sieht man, dass es auch nur endlich viele Werte für $Y$ gibt, an denen eine gemeinsame Nullstelle vorliegen kann. Damit kann es überhaupt nur endlich viele gemeinsame Nullstellen geben.


}


\bild{ \begin{center}


\includegraphics[width=5.5cm]{ \bildeinlesungpng {Two_cubic_curves} {png} }


\end{center}


\bildtext {} }


\bildlizenz { Two cubic curves.png } {} {Hack} {Commons} {PD} {}


\inputfaktbeweis


{Primpolynom/Ebene Kurve/Unendlich viele Punkte/Irreduzibel/Fakt}


{Korollar}


{}


{


\faktsituation {}


\faktvoraussetzung {Es sei $K$ ein
\definitionsverweis {Körper}{}{}
und sei


\mavergleichskette


{\vergleichskette


{ F
}


{ \in }{ K[X,Y]
}


{  }{
}


{    }{
}


{   }{
}


}
{}{}{}
ein
\definitionsverweis {Primpolynom}{}{.}
Die Kurve $V(F)$ besitze unendlich viele Elemente.}


\faktfolgerung {Dann ist das
\definitionsverweis {Verschwindungsideal}{}{}
zu $V(F)$ gleich dem
\definitionsverweis {Hauptideal}{}{}
$(F)$ und $V(F)$ ist
\definitionsverweis {irreduzibel}{}{.}}


\faktzusatz {}


\faktzusatz {}


}


{


Es ist


\mavergleichskette


{\vergleichskette


{ (F)
}


{ \subseteq }{ \operatorname{Id} \,(V(F))
}


{  }{
}


{    }{
}


{   }{
}


}
{}{}{.}
Es sei


\mavergleichskette


{\vergleichskette


{ G
}


{ \in }{ \operatorname{Id} \,(V(F))
}


{  }{
}


{    }{
}


{   }{
}


}
{}{}{.}
Nach

Lemma 3.8

 (3)

ist


\mavergleichskettedisp


{\vergleichskette


{ V(F)
}


{ =} { V (\operatorname{Id} \,(V(F)) )
}


{ \subseteq} { V(F,G)
}


{ } {
}


{ } {
}


}
{}{}{.}
Wenn $G$ kein Vielfaches von $F$ ist, so ergibt sich mit

Satz 4.8

direkt ein Widerspruch zur Voraussetzung. Also ist


\mavergleichskettedisp


{\vergleichskette


{ \operatorname{Id} \,(V(F))
}


{ =} {(F)
}


{ } {
}


{ } {
}


{ } {
}


}
{}{}{}
ein
\definitionsverweis {Primideal}{}{}
und nach

Lemma 4.3

ist $V(F)$ irreduzibel.


}
